\FloatBarrier
\section{Resultados}

\subsection{Sensibilidade ao tamanho do alvo}

A Figura~\ref{fig:sweep-pgf} mostra que o efeito do tamanho da pseudo-caixa foi
n\~ao monot\^onico. F1 e AP50 aumentaram das menores caixas at\'e uma regi\~ao de
estabilidade, mas o tamanho 160 apresentou queda acentuada e supercontagem. O
maior F1 foi 0,8253 em 192 pixels. O menor MAE (3,65) e o maior AP50 (0,8776)
ocorreram em 176 pixels. Em 144 pixels, F1 foi 0,8228, AP50 0,8748 e MAE 4,40.

\begin{figure}[htbp]
\centering
\begin{tikzpicture}
\begin{groupplot}[
 group style={group size=2 by 1,horizontal sep=1.15cm},
 width=.45\textwidth,height=4.65cm,
 xlabel={Lado da pseudo-caixa (px)},
 xmin=40,xmax=200,xtick={48,72,96,120,144,168,192},
 tick label style={font=\scriptsize},label style={font=\small},
 grid=major,grid style={black!10},axis line style={black!65}
]
\nextgroupplot[ylabel={Desempenho},ymin=.62,ymax=.96,legend style={font=\scriptsize,draw=none,at={(.5,-.36)},anchor=north,legend columns=3}]
\addplot[black,thick,mark=*] table[x=box_size_px,y=f1_val_iou50,col sep=comma]{dados/bbox_sweep.csv};
\addlegendentry{F1@0,50}
\addplot[black!50,thick,dashed,mark=square*] table[x=box_size_px,y=ap50_val,col sep=comma]{dados/bbox_sweep.csv};
\addlegendentry{AP50}
\addplot[black!30,thick,densely dotted,mark=triangle*] table[x=box_size_px,y=recall_val_iou50,col sep=comma]{dados/bbox_sweep.csv};
\addlegendentry{Revoca\c{c}\~ao}
\addplot[only marks,mark=*,mark size=3pt] coordinates {(144,0.8228176)};
\nextgroupplot[ylabel={MAE de contagem},ymin=0,ymax=29]
\addplot[black!55,thick,mark=triangle*] table[x=box_size_px,y=mae_contagem_val,col sep=comma]{dados/bbox_sweep.csv};
\addplot[only marks,mark=triangle*,mark size=3pt] coordinates {(144,4.4)};
\end{groupplot}
\end{tikzpicture}
\caption{Sensibilidade ao tamanho do alvo na valida\c{c}\~ao. O marcador maior
indica a configura\c{c}\~ao de 144 pixels usada na compara\c{c}\~ao de capacidade.}
\label{fig:sweep-pgf}
\end{figure}

Os valores entre 96 e 192 pixels, exceto 160, formam um plat\^o estreito de F1
(0,8157--0,8253); logo, os dados n\~ao sustentam um \'unico tamanho universal. O
ponto de 144 pixels foi mantido por combinar AP50 pr\'oxima do m\'aximo
(0,8748 vs.\ 0,8776 em 176), revoca\c{c}\~ao superior a 176 e 192
(0,8285 vs.\ 0,7885 e 0,8059) e posi\c{c}\~ao intermedi\'aria no plat\^o. A
instabilidade em 160 pixels n\~ao deve ser lida como benef\'icio de caixas maiores:
nesse ponto, o detector saturou o teto de 300 detec\c{c}\~oes por imagem na
valida\c{c}\~ao, com revoca\c{c}\~ao 0,955 e precis\~ao 0,496, sugerindo problema de
calibra\c{c}\~ao do ponto operacional em um treino isolado.

\subsection{Capacidade do modelo e limiar}

Com pseudo-caixas de 144 pixels, YOLO11s apresentou o maior F1 (0,8363), AP50
(0,8848) e o menor MAE (4,68) na valida\c{c}\~ao. YOLO11n ficou pr\'oximo, com F1
0,8231 e MAE 4,83. YOLO11m atingiu a maior revoca\c{c}\~ao (0,8921), mas sua
precis\~ao de 0,7158 elevou o MAE para 8,13 e reduziu F1 para 0,7943. A vantagem
do YOLO11s sobre YOLO11n foi pequena (0,0132 em F1 e 0,0075 em AP50), portanto
suporta a escolha operacional neste experimento, n\~ao uma afirma\c{c}\~ao geral de
superioridade arquitetural. O YOLO11m foi descartado porque o aumento de
revoca\c{c}\~ao n\~ao compensou falsos positivos e erro de contagem.

Esse comportamento ilustra por que mAP de treinamento n\~ao deve ser o \'unico
crit\'erio. Detectar mais candidatos pode ser desej\'avel em uma etapa de busca,
mas transfere custo para supress\~ao de duplicidades ou revis\~ao posterior,
sobretudo em campos densos. Assim, a escala do modelo n\~ao foi tratada como
proxy autom\'atica de qualidade.

Na sele\c{c}\~ao do ponto operacional, limiares baixos preservaram quase todas as
c\'elulas, mas geraram muitas detec\c{c}\~oes excedentes: em 0,10, a revoca\c{c}\~ao
foi 0,9661 e a precis\~ao 0,5187, com MAE 24,78. Em 0,50, a precis\~ao subiu para
0,8853, mas a revoca\c{c}\~ao caiu para 0,6919 e o vi\'es tornou-se negativo
($-6,28$ c\'elulas por imagem). O limiar 0,35 maximizou F1@0,50 (0,8363), com
precis\~ao 0,7983, revoca\c{c}\~ao 0,8782, MAE 4,68 e vi\'es 2,88. O teste recebeu
0,35 sem novo ajuste. A proximidade entre as curvas em IoU 0,30 e 0,50 indica
que, nesse intervalo, a maioria das associa\c{c}\~oes aceitas no crit\'erio mais
flex\'ivel tamb\'em satisfaz o crit\'erio de 0,50. A separa\c{c}\~ao em IoU 0,75
evidencia a penaliza\c{c}\~ao adicional por desalinhamento fino com a
pseudo-caixa; portanto, o resultado de teste mede transfer\^encia do ponto
operacional, n\~ao o melhor limiar retrospectivo para aquela parti\c{c}\~ao.

\subsection{Teste retido}

Com YOLO11s, lado 144 e confian\c{c}a 0,35, o teste produziu 975 verdadeiros
positivos, 257 falsos positivos e 174 falsos negativos. A
Tabela~\ref{tab:teste-unificado} re\'une as medidas com intervalo de 95\% obtido
por bootstrap de imagens.

\begin{table}[htbp]
\centering
\caption{Resultado no teste retido. IC95\% por bootstrap de imagens.}
\label{tab:teste-unificado}
\small
\setlength{\tabcolsep}{4.5pt}
\begin{tabular}{lrl@{\hspace{6mm}}lrl}
\toprule
M\'etrica & Valor & IC95\% & M\'etrica & Valor & IC95\% \\
\midrule
Precis\~ao & 0,7914 & [0,746; 0,833] & AP30 & 0,8843 & [0,858; 0,913] \\
Revoca\c{c}\~ao & 0,8486 & [0,822; 0,877] & AP50 & 0,8723 & [0,843; 0,904] \\
F1 & 0,8190 & [0,796; 0,842] & AP75 & 0,7515 & [0,695; 0,811] \\
MAE & 4,675 & [3,150; 6,425] & RMSE & 7,216 & [4,560; 9,850] \\
Vi\'es & 2,075 & [0,050; 4,126] & Correla\c{c}\~ao & 0,9434 & [0,874; 0,979] \\
\bottomrule
\end{tabular}
\end{table}

Em m\'edia, o teste teve 24,4 verdadeiros positivos, 6,4 falsos positivos e 4,4
falsos negativos por imagem. O F1 de teste ficou 0,0173 abaixo da valida\c{c}\~ao,
enquanto AP50 caiu 0,0125, diferen\c{c}as compat\'iveis com transfer\^encia
moderada entre duas amostras de 40 imagens. O vi\'es positivo indica 2,08
detec\c{c}\~oes excedentes por imagem em m\'edia. A Figura~\ref{fig:qualitativo}
mostra dois casos reais do teste: no primeiro, falsos positivos e falsos
negativos se compensam na contagem; no segundo, a revoca\c{c}\~ao completa vem
acompanhada de supercontagem por falsos positivos.

\begin{figure}[htbp]
\centering
\includegraphics[width=.90\textwidth]{qualitativo_teste.png}
\caption{Exemplos qualitativos do teste com YOLO11s, pseudo-caixa de 144 pixels
e confian\c{c}a 0,35. No cabe\c{c}alho, Anotadas \'e o total de c\'elulas de
refer\^encia e Preditas \'e o total de detec\c{c}\~oes. VP indica verdadeiro
positivo (caixa verde), FP falso positivo (caixa laranja com X) e FN falso
negativo (c\'irculo vermelho), no pareamento por IoU 0,50.}
\label{fig:qualitativo}
\end{figure}

Esses exemplos mostram por que a contagem n\~ao deve ser lida isoladamente. Uma
imagem pode ter erro absoluto nulo e ainda conter trocas locais entre omiss\~oes e
detec\c{c}\~oes excedentes; outra pode recuperar todos os pontos anotados e,
mesmo assim, gerar carga de revis\~ao por duplicidades ou estruturas confundidas
com c\'elulas. Por isso, AP, F1, MAE e inspe\c{c}\~ao qualitativa s\~ao
complementares neste problema.

Na an\'alise explorat\'oria por tercis, o MAE foi 4,15 em imagens com 5--16
c\'elulas, 3,69 entre 17--30 e 6,07 entre 32--94. A faixa mais densa concentrou
maior erro absoluto, embora o F1 global permane\c{c}a 0,8190. A mediana global do
erro foi 3 c\'elulas. Esses estratos foram definidos depois da avalia\c{c}\~ao
principal e n\~ao devem ser interpretados como hip\'oteses confirmat\'orias, mas
ajudam a localizar onde a contagem se torna mais sens\'ivel a duplicidades e
omiss\~oes.

A avalia\c{c}\~ao por dist\^ancia ao ponto confirma que o resultado n\~ao depende
apenas da geometria da pseudo-caixa. Como n\'ucleos vizinhos distam em mediana 72
pixels, raios abaixo desse valor associam cada predi\c{c}\~ao ao n\'ucleo correto,
n\~ao a um adjacente. No teste, o F1 foi 0,812 em $R=32$ pixels e 0,821 em
$R=48$, coincidindo com o F1 por IoU 0,50 (0,819). Duas medidas de natureza
distinta, portanto, indicam localiza\c{c}\~ao efetiva dos n\'ucleos anotados.

\subsection{Interpreta\c{c}\~ao}

O experimento n\~ao recupera fronteiras celulares; ele aprende uma vizinhan\c{c}a
quadrada ao redor do n\'ucleo. AP75 alta indica concord\^ancia com a regra da
Equa\c{c}\~ao~\ref{eq:pseudobox}, n\~ao segmenta\c{c}\~ao precisa de n\'ucleo ou
citoplasma. Essa distin\c{c}\~ao tamb\'em limita compara\c{c}\~oes diretas com
m\'etodos de segmenta\c{c}\~ao ou classifica\c{c}\~ao Bethesda, pois m\'etrica,
unidade de an\'alise, parti\c{c}\~ao e sem\^antica do r\'otulo precisam coincidir
para sustentar superioridade num\'erica
\cite{jiang2023systematic,kalbhor2023cervicell}.

A obje\c{c}\~ao de que o IoU contra um alvo sint\'etico seria igualmente
artificial \'e respondida pela avalia\c{c}\~ao por dist\^ancia: como o F1 por ponto
converge ao F1 por IoU, a m\'etrica n\~ao depende criticamente da geometria
adotada, e o n\'ucleo anotado -- \'unico ground truth espacial da CRIC -- \'e de
fato recuperado. Ainda assim, pontos cont\^em informa\c{c}\~ao de inst\^ancia que
r\'otulos globais n\~ao possuem. A literatura de supervis\~ao pontual mostra que
essa informa\c{c}\~ao pode orientar localiza\c{c}\~ao e segmenta\c{c}\~ao com menor
custo de anota\c{c}\~ao \cite{chen2021sparsepoints,xia2023colorization}.
O resultado da RQ1 acrescenta uma observa\c{c}\~ao pr\'atica: quando pontos s\~ao
convertidos em caixas, a geometria escolhida deve ser reportada e testada.

A correla\c{c}\~ao 0,9434 mostra que o modelo acompanha a varia\c{c}\~ao de carga
celular entre imagens; o MAE de 4,68 e o vi\'es de 2,08 revelam erro residual e
tend\^encia de supercontagem. As causas prov\'aveis, n\~ao separadas por este
experimento, s\~ao duplicidade em grupos sobrepostos, estruturas n\~ao celulares
confundidas com o alvo e c\'elulas vis\'iveis sem ponto correspondente na
anota\c{c}\~ao. Distinguir essas causas exige revis\~ao cega de falsos positivos por
especialista.
